\documentclass[12pt, a4paper]{article}

\usepackage[T1]{fontenc}
\usepackage[utf8]{inputenc}
\usepackage[french]{babel}

\usepackage{amsmath}
\usepackage{amssymb}
\usepackage{graphicx}
\usepackage{circuitikz}

\title{Préparation de TP : L'Amplificateur Opérationnel}
\date{}



\begin{document}
\maketitle

\tableofcontents 

\newpage

\section{Introduction}
Ce TP porte sur l'étude expérimentale de l'Amplificateur Opérationnel (AOP) et du circuit intégré NE555. Nous allons explorer leurs limites (Slew-Rate) et leurs applications (filtrage, oscillation, modulation).



\section{Mesure du Slew-Rate}
\subsection{Définition théorique}
Le Slew-Rate (SR) représente la vitesse maximale de variation de la tension de sortie de l'AOP (unité Volt/µs). 

\subsection{Calculs préparatoires}
Le montage utilisé est un montage non-inverseur de gain 1. En appliquant un signal carré de forte amplitude en entrée, la sortie ne peut pas suivre instantanément et présente une rampe. 
Le SR sera calculé graphiquement sur l'oscilloscope par :
$$SR = \frac{\Delta V_s}{\Delta t} \text{ (en } V/\mu s\text{)}$$


\section{Montage Pseudo-Intégrateur}
\subsection{Détermination de la fonction de transfert}
D'après le schéma, en utilisant le théorème de Millman à l'entrée inverseuse (en supposant l'AOP idéal en régime linéaire) :
$$H(j\omega) = \frac{\underline{V_s}}{\underline{V_e}} = -\frac{Z_2}{R_1}$$
Avec $Z_2 = \frac{R_2 \cdot \frac{1}{jC\omega}}{R_2 + \frac{1}{jC\omega}} = \frac{R_2}{1 + jR_2C\omega}$.
D'où :
$$H(j\omega) = -\frac{R_2}{R_1} \cdot \frac{1}{1 + jR_2C\omega}$$

\subsection{Fréquence de coupure et condition d'intégration}
La fonction d'intégration est réalisée pour des fréquences $f \gg f_c$ avec $f_c = \frac{1}{2\pi R_2 C}$.
Application numérique : $R_2 = 100~k\Omega$, $C = 10~nF$.
$$f_c = \frac{1}{2\pi \cdot 10^5 \cdot 10^{-8}} \approx 159~Hz$$


\section{Oscillateur sinusoïdal à pont de Wien}
\subsection{Condition d'oscillation}
Le gain de l'amplificateur non-inverseur est $A = 1 + \frac{R_2}{R_1}$. Pour compenser les pertes du pont de Wien, il faut $A \geq 3$ ($\beta = \frac{1}{3}$).
On doit donc régler le potentiomètre $R_1$ tel que $R_1 \leq \frac{R_2}{2}$.

\subsection{Fréquence propre}
La fréquence des oscillations est donnée par :
$$f_0 = \frac{1}{2\pi RC} = \frac{1}{2\pi \cdot 10^3 \cdot 22 \cdot 10^{-9}} \approx 7234~Hz$$

\section{Génération de signaux avec le NE555}
\subsection{Montage Astable}
La fréquence du signal de sortie est $f = \frac{1,44}{(R_a + 2R_b)C}$.
Pour obtenir $f \approx 700~Hz$ avec $R_a = 1~k\Omega$ et $C = 0,1~\mu F$:
$$R_b = \frac{1}{2} \left( \frac{1,44}{f \cdot C} - R_a \right) = \frac{1}{2} \left( \frac{1,44}{700 \cdot 10^{-7}} - 1000 \right) \approx 9,78~k\Omega$$


https://www.electronics-tutorials.ws/waveforms/555-circuits-part-1.html

\end{document}